% Môn: Vật lý
% Lớp: 11
% Chương: Chương I. Dao động
% Bài: L11C1 Bài 1. Dao động điều hòa · Xác định các đại lượng đặc trưng dựa vào công thức (từ phương trình dao động).
% ===== Câu 36 =====
% ID: L11C1B1-152-DS
% Mức: TH
\dangbt{Xác định các đại lượng đặc trưng dựa vào công thức (từ phương trình dao động).}
\begin{ex}
% ID: L11C1B1-152-DS
% Mức: TH
Bài 4 – xác định đại lượng từ phương trình: Với $x=5\cos(4\pi t-\pi/6)$ cm.
\choiceTF
{\True Biên độ bằng $5\,\text{cm}$.}
{\True Tần số bằng $2\,\text{Hz}$.}
{\True Chu kì bằng $0,5\,\text{s}$.}
{Pha ban đầu bằng $\pi/6$ rad.}
\end{ex}

% ===== Câu 37 =====
% ID: L11C1B1-153-DS
% Mức: TH
\begin{ex}
% ID: L11C1B1-153-DS
% Mức: TH
Một vật dao động theo phương trình $x=4\cos(2\pi t+\pi/3)$ cm. Xác định biên độ, tần số góc, chu kì, tần số và pha ban đầu.
\choiceTF

\end{ex}

% ===== Câu 40 =====
% ID: L11C1B1-154-DS
% Mức: VD
\begin{ex}
% ID: L11C1B1-154-DS
% Mức: VD
Bài 4 – xác định đại lượng từ phương trình: Với $x=8\sin(2\pi t+\pi/3)$ cm.
\choiceTF
{\True Biên độ bằng $8\,\text{cm}$.}
{\True Tần số góc bằng $2\pi$ rad/s.}
{\True Chu kì bằng $1\,\text{s}$.}
{Tần số bằng $2\pi$ Hz.}
\end{ex}

% ===== Câu 41 =====
% ID: L11C1B1-155-DS
% Mức: VD
\begin{ex}
% ID: L11C1B1-155-DS
% Mức: VD
Một vật dao động theo phương trình $x=6\cos(4\pi t+\pi/3)$ cm. Xác định biên độ, tần số góc, chu kì, tần số và pha ban đầu.
\choiceTF

\end{ex}

% ===== Câu 44 =====
% ID: L11C1B1-156-DS
% Mức: VD
\begin{ex}
% ID: L11C1B1-156-DS
% Mức: VD
Bài 4 – xác định đại lượng từ phương trình: Với $x=6\cos(\pi t+\pi/2)$ cm.
\choiceTF
{\True Biên độ bằng $6\,\text{cm}$.}
{\True Tần số bằng $0,5\,\text{Hz}$.}
{\True Chu kì bằng $2\,\text{s}$.}
{Tại $t=0$, vật ở biên dương.}
\end{ex}

% ===== Câu 45 =====
% ID: L11C1B1-157-DS
% Mức: VD
\begin{ex}
% ID: L11C1B1-157-DS
% Mức: VD
Một vật dao động theo phương trình $x=8\cos(5\pi t+\pi/3)$ cm. Xác định biên độ, tần số góc, chu kì, tần số và pha ban đầu.
\choiceTF

\end{ex}

% ===== Câu 52 =====
% ID: L11C1B1-158-TLN
% Mức: TH
\begin{ex}
% ID: L11C1B1-158-TLN
% Mức: TH
Một vật dao động theo phương trình x = 6cos(4πt + π/3) cm. Chu kì dao động theo đơn vị s bằng bao nhiêu?
\shortans{0,5}
\loigiai{
$T$ = 2π/ω = 2π/(4π) = $0,5\,\text{s}$.
}
\end{ex}

% ===== Câu 54 =====
% ID: L11C1B1-159-TLN
% Mức: TH
\begin{ex}
% ID: L11C1B1-159-TLN
% Mức: TH
Một vật dao động theo phương trình x = 6cos(4πt + π/3) cm. Chu kì dao động theo đơn vị s bằng bao nhiêu?
\shortans{0,5}
\loigiai{
$T$ = 2π/ω = 2π/(4π) = $0,5\,\text{s}$.
}
\end{ex}

% ===== Câu 56 =====
% ID: L11C1B1-160-TLN
% Mức: TH
\begin{ex}
% ID: L11C1B1-160-TLN
% Mức: TH
Một vật dao động theo phương trình x = 6cos(4πt + π/3) cm. Chu kì dao động theo đơn vị s bằng bao nhiêu?
\shortans{0,5}
\loigiai{
$T$ = 2π/ω = 2π/(4π) = $0,5\,\text{s}$.
}
\end{ex}

% ===== Câu 58 =====
% ID: L11C1B1-161-TLN
% Mức: TH
\begin{ex}
% ID: L11C1B1-161-TLN
% Mức: TH
Một vật dao động theo phương trình x = 8cos(5πt - π/4) cm. Tần số dao động theo đơn vị Hz bằng bao nhiêu?
\shortans{2,5}
\loigiai{
f = ω/(2π) = 5π/(2π) = $2,5\,\text{Hz}$.
}
\end{ex}

% ===== Câu 60 =====
% ID: L11C1B1-162-TLN
% Mức: TH
\begin{ex}
% ID: L11C1B1-162-TLN
% Mức: TH
Một vật dao động theo phương trình x = 8cos(5πt - π/4) cm. Tần số dao động theo đơn vị Hz bằng bao nhiêu?
\shortans{2,5}
\loigiai{
f = ω/(2π) = 5π/(2π) = $2,5\,\text{Hz}$.
}
\end{ex}

% ===== Câu 62 =====
% ID: L11C1B1-163-TLN
% Mức: TH
\begin{ex}
% ID: L11C1B1-163-TLN
% Mức: TH
Một vật dao động theo phương trình x = 8cos(5πt - π/4) cm. Tần số dao động theo đơn vị Hz bằng bao nhiêu?
\shortans{2,5}
\loigiai{
f = ω/(2π) = 5π/(2π) = $2,5\,\text{Hz}$.
}
\end{ex}

% ===== Câu 64 =====
% ID: L11C1B1-164-TN
% Mức: TH
\begin{ex}
% ID: L11C1B1-164-TN
% Mức: TH
Gia tốc của vật dao động điều hoà bằng 0 khi
\choice
{\True hợp lực tác dụng vào vật bằng 0}
{không có vị trí nào có gia tốc bằng 0}
{vật ở hai biên}
{vật ở vị trí có vận tốc bằng 0}
\loigiai{
Gia tốc tỉ lệ với hợp lực theo định luật II Newton: $a = \frac{F}{m}$. Gia tốc bằng 0 thì lực tổng hợp bằng 0.
}
\end{ex}

% ===== Câu 65 =====
% ID: L11C1B1-165-TN
% Mức: TH
\begin{ex}
% ID: L11C1B1-165-TN
% Mức: TH
Gia tốc của vật dao động điều hoà bằng 0 khi
\choice
{\True hợp lực tác dụng vào vật bằng 0}
{không có vị trí nào có gia tốc bằng 0}
{vật ở hai biên}
{vật ở vị trí có vận tốc bằng 0}
\loigiai{
Gia tốc tỉ lệ với hợp lực theo định luật II Newton: $a = \frac{F}{m}$. Gia tốc bằng 0 thì lực tổng hợp bằng 0.
}
\end{ex}

% ===== Câu 66 =====
% ID: L11C1B1-166-TN
% Mức: TH
\begin{ex}
% ID: L11C1B1-166-TN
% Mức: TH
Gia tốc của vật dao động điều hoà bằng 0 khi
\choice
{\True hợp lực tác dụng vào vật bằng 0}
{không có vị trí nào có gia tốc bằng 0}
{vật ở hai biên}
{vật ở vị trí có vận tốc bằng 0}
\loigiai{
Gia tốc tỉ lệ với hợp lực theo định luật II Newton: $a = \frac{F}{m}$. Gia tốc bằng 0 thì lực tổng hợp bằng 0.
}
\end{ex}

% ===== Câu 67 =====
% ID: L11C1B1-167-TN
% Mức: TH
\begin{ex}
% ID: L11C1B1-167-TN
% Mức: TH
Cho một chất điểm dao động điều hòa với phương trình: $x=3\sin(\omega t-5\pi/6)$ (cm). Pha ban đầu của dao động nhận giá trị nào sau đây
\choice
{\True $2\pi/3$ rad}
{$4\pi/3$ rad}
{$-5\pi/6$ rad}
{$\pi/3$ rad}
\loigiai{
$\sin(\omega t - 5\pi/6) = \cos(\omega t - 4\pi/3)$ nên pha ban đầu: $-4\pi/3 \equiv 2\pi/3 \pmod{2\pi}$.
}
\end{ex}

% ===== Câu 68 =====
% ID: L11C1B1-168-TN
% Mức: TH
\begin{ex}
% ID: L11C1B1-168-TN
% Mức: TH
Trong phương trình dao động điều hòa $x = A\cos(\omega t + \varphi)$, phát biểu nào sau đây sai?
\choice
{biên độ $A$ là hằng số dương, phụ thuộc vào kích thích dao động}
{biên độ $A$ là hằng số dương, không phụ thuộc vào gốc thời gian}
{\True pha ban đầu $\varphi$ là hằng số, chỉ phụ thuộc vào gốc thời gian}
{tần số góc $\omega$ là hằng số dương, phụ thuộc vào các đặc tính của hệ}
\loigiai{
Pha ban đầu $\varphi$ không chỉ phụ thuộc vào gốc thời gian mà còn vào điều kiện ban đầu của hệ thống.
}
\end{ex}

% ===== Câu 69 =====
% ID: L11C1B1-169-TN
% Mức: VD
\begin{ex}
% ID: L11C1B1-169-TN
% Mức: VD
Một con lắc lò xo gồm vật nặng có khối lượng $0{,}2$ kg và lò xo có độ cứng $80$ N/m. Con lắc lò xo dao động điều hòa với biên độ $3$ cm. Tốc độ cực đại của vật nặng bằng:
\choice
{\True $0{,}6$ m/s}
{$0{,}7$ m/s}
{$0{,}5$ m/s}
{$0{,}4$ m/s}
\loigiai{
$\omega=\sqrt{\dfrac{k}{m}}=\sqrt{\dfrac{80}{0,2}}=20$,
$v_\text{max}=\omega A=20\cdot 0,03=0,6$ m/s.
}
\end{ex}

% ===== Câu 70 =====
% ID: L11C1B1-170-TN
% Mức: VD
\begin{ex}
% ID: L11C1B1-170-TN
% Mức: VD
Một vật dao động điều hòa với phương trình $x=-5\cos(4\pi t)(\mathrm{~cm})$. Biên độ và pha ban đầu của dao động lần lượt là
\choice
{$5$ cm; $0$ rad}
{$-5$ cm; $4\pi$ rad}
{$5$ cm; $4\pi$ rad}
{\True $5$ cm; $\pi$ rad}
\loigiai{
Phương trình có dạng $x=A\cos(\omega$ t $+\varphi)$.
Một vật dao động điều hòa với phương trình $x=-5\cos(4\pi t)(\mathrm{~cm})=5\cos(4\pi t+\pi)$
Ta có biên độ $A=5\mathrm{~cm}$ và pha ban đầu $\varphi=\pi$ rad
}
\end{ex}

% ===== Câu 71 =====
% ID: L11C1B1-171-TN
% Mức: VD
\begin{ex}
% ID: L11C1B1-171-TN
% Mức: VD
Một vật dao động điều hòa với phương trình $x=-5\cos(4\pi t)(\mathrm{~cm})$. Biên độ và pha ban đầu của dao động lần lượt là
\choice
{$5$ cm; $0$ rad}
{$-5$ cm; $4\pi$ rad}
{$5$ cm; $4\pi$ rad}
{\True $5$ cm; $\pi$ rad}
\loigiai{
Phương trình có dạng $x=A\cos(\omega$ t $+\varphi)$.
Một vật dao động điều hòa với phương trình $x=-5\cos(4\pi t)(\mathrm{~cm})=5\cos(4\pi t+\pi)$
Ta có biên độ $A=5\mathrm{~cm}$ và pha ban đầu $\varphi=\pi$ rad
}
\end{ex}

% ===== Câu 72 =====
% ID: L11C1B1-172-TN
% Mức: VD
\begin{ex}
% ID: L11C1B1-172-TN
% Mức: VD
Một vật dao động điều hòa với phương trình $x=-5\cos(4\pi t)(\mathrm{~cm})$. Biên độ và pha ban đầu của dao động lần lượt là
\choice
{$5$ cm; $0$ rad}
{$-5$ cm; $4\pi$ rad}
{$5$ cm; $4\pi$ rad}
{\True $5$ cm; $\pi$ rad}
\loigiai{
Phương trình có dạng $x=A\cos(\omega$ t $+\varphi)$.
Một vật dao động điều hòa với phương trình $x=-5\cos(4\pi t)(\mathrm{~cm})=5\cos(4\pi t+\pi)$
Ta có biên độ $A=5\mathrm{~cm}$ và pha ban đầu $\varphi=\pi$ rad
}
\end{ex}

% ===== Câu 73 =====
% ID: L11C1B1-173-TN
% Mức: VD
\begin{ex}
% ID: L11C1B1-173-TN
% Mức: VD
Một vật dao động điều hoà theo phương trình $x=8\cos(2\pi t - \pi/2)$ cm. Vận tốc và gia tốc của vật khi vật đi qua ly độ $4\sqrt{3}$ cm là
\choice
{$-8\pi$ cm/s và $16\pi^2\sqrt{3}$ cm/s $^2$}
{$8\pi$ cm/s và $16\pi^2$ cm/s $^2$}
{$\pm 8\pi$ cm/s và $\pm 16\pi^2\sqrt{3}$ cm/s $^2$}
{\True $\pm 8\pi$ cm/s và $-6\pi^2\sqrt{3}$ cm/s $^2$}
\loigiai{
$\omega = 2\pi$, $A = 8$. Với $x = 4\sqrt{3}$,
$v = \omega \sqrt{A^2 - x^2} = 2\pi \cdot \sqrt{64 - 48} = 2\pi \cdot 4 = 8\pi$ (cm/s).
$a = -\omega^2 x = -4\pi^2 \cdot 4\sqrt{3} = -16\pi^2\sqrt{3}$.
}
\end{ex}

% ===== Câu 74 =====
% ID: L11C1B1-174-TN
% Mức: VD
\begin{ex}
% ID: L11C1B1-174-TN
% Mức: VD
Một vật dao động điều hoà theo phương trình $x=8\cos(2\pi t - \pi/2)$ cm. Vận tốc và gia tốc của vật khi vật đi qua ly độ $4\sqrt{3}$ cm là
\choice
{$-8\pi$ cm/s và $16\pi^2\sqrt{3}$ cm/s $^2$}
{$8\pi$ cm/s và $16\pi^2$ cm/s $^2$}
{$\pm 8\pi$ cm/s và $\pm 16\pi^2\sqrt{3}$ cm/s $^2$}
{\True $\pm 8\pi$ cm/s và $-6\pi^2\sqrt{3}$ cm/s $^2$}
\loigiai{
$\omega = 2\pi$, $A = 8$. Với $x = 4\sqrt{3}$,
$v = \omega \sqrt{A^2 - x^2} = 2\pi \cdot \sqrt{64 - 48} = 2\pi \cdot 4 = 8\pi$ (cm/s).
$a = -\omega^2 x = -4\pi^2 \cdot 4\sqrt{3} = -16\pi^2\sqrt{3}$.
}
\end{ex}

% ===== Câu 75 =====
% ID: L11C1B1-175-TN
% Mức: VD
\begin{ex}
% ID: L11C1B1-175-TN
% Mức: VD
Một vật dao động điều hoà theo phương trình $x=8\cos(2\pi t - \pi/2)$ cm. Vận tốc và gia tốc của vật khi vật đi qua ly độ $4\sqrt{3}$ cm là
\choice
{$-8\pi$ cm/s và $16\pi^2\sqrt{3}$ cm/s $^2$}
{$8\pi$ cm/s và $16\pi^2$ cm/s $^2$}
{$\pm 8\pi$ cm/s và $\pm 16\pi^2\sqrt{3}$ cm/s $^2$}
{\True $\pm 8\pi$ cm/s và $-6\pi^2\sqrt{3}$ cm/s $^2$}
\loigiai{
$\omega = 2\pi$, $A = 8$. Với $x = 4\sqrt{3}$,
$v = \omega \sqrt{A^2 - x^2} = 2\pi \cdot \sqrt{64 - 48} = 2\pi \cdot 4 = 8\pi$ (cm/s).
$a = -\omega^2 x = -4\pi^2 \cdot 4\sqrt{3} = -16\pi^2\sqrt{3}$.
}
\end{ex}
